\section{The distributed point function}
\label{sec:tech:dpf}

A point function $f_{\alpha,\beta}$ takes the value $\beta$ at $x = \alpha$ and
zero everywhere else. A $(2,2)$-distributed point function splits it into two
keys, each of which individually hides both $\alpha$ and $\beta$, such that the
two evaluations recombine to $f_{\alpha,\beta}(x)$. The construction is that of
Boyle, Gilboa and Ishai, and it is written by us rather than imported, because
it is the component a viva will probe hardest and because the same
implementation serves both halves of this project: the private read layer in
Section~\ref{sec:tech:pir}, and the function secret sharing gates that private
training will need in the next phase.

\subsection{Why the keys are small}
\label{sec:tech:dpf:size}

The construction is a binary tree of pseudorandom seeds. Each key holds one
initial seed and one correction word per level, so key size is
\emph{logarithmic} in the domain while a naive one-hot query is linear. Table
\ref{tab:tech:keysize} gives measured sizes from the serialiser rather than
from the closed form, since the closed form is what the design draft got wrong.

\begin{table}[h]
\centering
\small
\caption{Measured DPF key size against the naive alternative. The packed
format is $1 + 4 + 16 + 18\,\mathrm{db} + \lceil \bits/8 \rceil$ bytes, and
\texttt{SizeBytes()} is asserted equal to the actual serialised length.}
\label{tab:tech:keysize}
\begin{tabular}{@{}rrrrrr@{}}
\toprule
domain bits & $\items$ & key, $\bits=64$ & key, $\bits=128$ & naive one-hot & compression \\
\midrule
10 & 1{,}024  & 209 B & 217 B & 8{,}192 B   & $39\times$ \\
11 & 2{,}048  & 227 B & 235 B & 16{,}384 B  & $72\times$ \\
12 & 4{,}096  & 245 B & 253 B & 32{,}768 B  & $134\times$ \\
16 & 65{,}536 & 317 B & 325 B & 524{,}288 B & $1654\times$ \\
\bottomrule
\end{tabular}
\end{table}

The design draft claims roughly 260 bytes at $\items = 4096$ against 32~KB
naive. The 32~KB figure is right. The 260-byte figure is not: the measured
value is \textbf{245 bytes}, and the discrepancy comes from the draft's formula
omitting both the initial seed and the self-describing header, while
simultaneously evaluating at 16 domain bits rather than the 12 that
$\items = 4096$ implies. The number in this report is measured by a test that
prints it, which is why the error surfaced.

\subsection{The sign convention}
\label{sec:tech:dpf:sign}

Textbook presentations of this construction define
$\Eval_b = (-1)^b\bigl(\mathrm{Convert}(s) + t \cdot \mathrm{CW}_{\text{last}}\bigr)$,
which yields the \emph{sum} convention $\Eval_0(x) + \Eval_1(x) = f(x)$. The
specification for this project instead states the invariant as a
\emph{difference}. We therefore omit the $(-1)^b$ factor and return
$\mathrm{Convert}(s) + t \cdot \mathrm{CW}_{\text{last}}$ directly, giving
%
\[
  \Eval(k_0, x) - \Eval(k_1, x) =
  \begin{cases}
    \beta & x = \alpha \\
    0     & \text{otherwise.}
  \end{cases}
\]
%
The two conventions are algebraically identical and differ only in where the
negation is placed. The reason to record the choice explicitly is that getting
it backwards makes every test fail in a way that resembles a tree traversal
bug rather than a sign error, which is an expensive place to spend a day.

\subsection{Testing: the structural invariant, checked at every node}
\label{sec:tech:dpf:testing}

A black-box test reports only that a leaf is wrong. The characteristic failure
of this construction is a mistaken control-bit update, which stays invisible
until enough levels have accumulated for it to change an output, and therefore
presents as ``correct at two domain bits, wrong at eight''. Bisecting that from
leaf values is slow.

The test suite therefore walks \emph{both} parties' trees in lockstep and
asserts the underlying structural invariant at every internal node. Writing
$(s_b, t_b)$ for party $b$'s state:

\begin{center}
\begin{tabular}{@{}ll@{}}
\toprule
On the path to $\alpha$  & $s_0 \neq s_1$ and $t_0 \oplus t_1 = 1$ \\
Off the path             & $s_0 = s_1$ and $t_0 \oplus t_1 = 0$ \\
\bottomrule
\end{tabular}
\end{center}

Off the path the two parties hold identical state, which is precisely why
their contributions cancel. On the path they differ and their control bits
disagree, which is what lets the final correction word inject $\beta$ at
exactly one leaf. Every correctness property of the scheme follows from these
two lines, so a violation names the exact level and node at which the
construction first went wrong.

This checker was written \emph{before} the generator was trusted, and both
were exercised together on the first run.

\subsection{Correctness results}
\label{sec:tech:dpf:results}

The structural invariant holds at every node for every $\alpha$ in every
domain up to eight bits, and at sampled $\alpha$ up to eleven bits.

The difference invariant is checked against a brute-force point function
oracle for every $\alpha$ and every $x$: exhaustively to ten domain bits on
the 64-bit ring and nine on the 128-bit ring, then for all $x$ at sampled
$\alpha$ at eleven and twelve bits. Edge cases that off-by-one errors land on
are covered explicitly, namely $\alpha$ at the first and last index, $\beta = 0$
where everything must cancel, and $\beta = -1$ where the ring wraps. A separate
check confirms that $\beta = 1$ reconstructs exactly, since the read layer
depends on that with no tolerance for an error term.

Traversal is implemented once and shared by single-point evaluation,
whole-domain evaluation and the test checker, so there is one traversal
implementation to be wrong rather than three.

\subsection{What is implemented but not yet verified}
\label{sec:tech:dpf:pending}

Whole-domain evaluation and key serialisation are implemented, because the
header is a frozen cross-component contract and had to be complete, but they
are so far exercised only indirectly. Extending the exhaustive invariant to
sixteen domain bits through whole-domain evaluation, asserting that it agrees
with single-point evaluation at every index, and round-tripping serialised keys
are the next steps. The key sizes in
Table~\ref{tab:tech:keysize} do come from the real serialiser, so that path is
at least exercised.
